\documentclass[11pt]{scrartcl}

\usepackage[top=1.5cm]{geometry}
\usepackage{url}
\usepackage{float}
\usepackage{listings}
\usepackage{xcolor}

\setlength{\parindent}{0em}
\setlength{\parskip}{0.5em}

\newcommand{\youranswerhere}{[Your answer goes here \ldots]}
\renewcommand{\thesubsection}{\arabic{subsection}}

\lstdefinestyle{dbtsql}{
  language=SQL,
  basicstyle=\small\ttfamily,
  keywordstyle=\color{magenta!75!black},
  stringstyle=\color{green!50!black},
  showspaces=false,
  showstringspaces=false,
  commentstyle=\color{gray}}

\title{
  \textbf{\large Assignment 2} \\
  Query Tuning \\
  {\large Database Tuning}
}

\author{
  Group Name (B15) \\
  \large Kitajima Mitsuhiro, 12314106 \\
  \large Sadykova Nargiz, 12129337 \\
  \large El Lbidi Omar, 12046411 \\
}

\begin{document}

\maketitle\thispagestyle{empty}

\subsection*{Creating Tables and Indexes}

SQL statements used to create the tables \texttt{Employee}, \texttt{Student}, and \texttt{Techdept}, and the indexes on the tables:

\begin{lstlisting}[style=dbtsql]
CREATE TABLE employee (
  ssnum INT, 
  name TEXT, 
  manager TEXT, 
  dept TEXT, 
  salary INT, 
  numfriends INT
);
CREATE UNIQUE INDEX idx_ssnum_emp ON employee (ssnum);
CREATE UNIQUE INDEX idx_name_emp ON employee (name);
CREATE INDEX idx_dept_emp ON employee (dept);

CREATE TABLE student (
  ssnum INT,
  name TEXT,
  course TEXT,
  grade FLOAT
);
CREATE UNIQUE INDEX idx_ssnum_stud ON student (ssnum);
CREATE UNIQUE INDEX idx_name_stud ON student (name);

CREATE TABLE techdept (
  dept TEXT,
  manager TEXT,
  location TEXT
);
CREATE UNIQUE INDEX idx_dept_tech ON techdept (dept);
\end{lstlisting}

\subsection*{Populating the Tables}

How did you fill the tables? What values did you use? Give a short description of your program.

We used a Python script (\texttt{populate.py}) utilizing the \texttt{pandas} and \texttt{random} libraries to generate the data and write it to TSV files. We generated 100,000 employees, 100,000 students, and 10 technical departments. To avoid key collisions, \texttt{ssnum} was generated by adding an offset to the loop index: 100,000 for employees and 200,000 for students. 

For the \texttt{employee} table, \texttt{name} was generated as ``Employee'' concatenated with the index. \texttt{manager} was randomly chosen from a list of 7 managers. About 10\% of employees were randomly assigned to one of 10 technical departments (``Tech1'' to ``Tech10''), while the remaining 90\% were assigned to one of 10 non-technical departments. \texttt{salary} was randomly generated between 50,000 and 1,500,000, and \texttt{numfriends} between 0 and 20.

For the \texttt{student} table, \texttt{name} was generated as ``Student'' concatenated with the index. \texttt{course} was randomly selected from a list of 10 courses, and \texttt{grade} was randomly generated between 1.0 and 5.0.

For the \texttt{techdept} table, each of the 10 technical departments was assigned a \texttt{manager} randomly selected from the 7 managers and a \texttt{location} randomly selected from 5 locations.

\subsection*{Queries}

\subsubsection*{Query 1}

\paragraph{Original Query}

Give the first type of query that might be hard for your database to optimize.

\begin{lstlisting}[style=dbtsql]
SELECT E.ssnum 
FROM employee E 
WHERE EXISTS 
  (SELECT 1 
  FROM techdept T 
  WHERE T.dept = E.dept)
\end{lstlisting}

\paragraph{Rewritten Query}

Give the rewritten query.

\begin{lstlisting}[style=dbtsql]
SELECT E.ssnum 
FROM employee E 
JOIN techdept T 
ON T.dept = E.dept
\end{lstlisting}

\paragraph{Evaluation of the Execution Plans}

Give the execution plan of the original query.

\textbf{PostgreSQL:}

{\small
\parskip0pt\begin{verbatim}
  Planning time:   4.603 ms
  Execution time:  39.585 ms

  Hash Join  [Inner]
    Actual time:   startup=0.196 ms  total=39.224 ms
    Actual rows:   9891.0  |  Hash Cond: (e.dept = t.dept)

    +-  Seq Scan on employee as e  (Outer)
        Actual time:   total=32.928 ms  |  Actual rows: 100000

    +-  Hash  (Inner)  Peak memory: 9 KB
        Actual time:   startup=0.185 ms
        Actual rows:   10.0

        +-  Seq Scan on techdept as t
            Actual rows:   10.0
\end{verbatim}}

\textbf{MariaDB:}

{\small
\parskip0pt\begin{verbatim}
  Optimization time: 10.5053 ms
  Total time:        66.7112 ms

  Duplicates Removal
    +-  Table: T  [ALL]
        Key: none  |  Rows: 10  |  pages_accessed=1
        Time: table=1.6580 ms

    +-  Table: E  [ref]  Key: idx_dept_emp
        Ref: mariadb.T.dept
        Loops: planned=10  actual=10
        Rows:  estimated=1  actual=989.1  |  pages_accessed=19829
        Time:  table=58.3611 ms
\end{verbatim}}

Give an interpretation of the execution plan, i.e., describe how the original query is evaluated.

\textbf{PostgreSQL} rewrites the correlated \texttt{EXISTS} subquery into an inner hash join. A sequential scan reads all 10 \texttt{techdept} rows and builds a small hash table (9\,KB, 0.185\,ms startup). A sequential scan then reads all 100,000 \texttt{employee} rows (32.928\,ms), probing the hash table on \texttt{dept} and returning 9,891 matches. The elevated planning time (4.603\,ms) reflects the cost of unnesting the correlated subquery.

\textbf{MariaDB} does not fully unnest the subquery. Instead it applies a \textit{Duplicates Removal} strategy: a full scan of \texttt{techdept} (10 rows, 1 page) produces 10 distinct \texttt{dept} values, and for each one the index \texttt{idx\_dept\_emp} is used to look up matching \texttt{employee} rows (10 loops, $\sim$989 rows each). Although the index avoids a full \texttt{employee} scan, the 10 batched lookups access 19,829 pages in total (58.361\,ms). The high optimisation time (10.505\,ms) reflects the cost of analysing the \texttt{EXISTS} form.

Give the execution plan of the rewritten query.

\textbf{PostgreSQL:}

{\small
\parskip0pt\begin{verbatim}
  Planning time:   0.093 ms
  Execution time:  9.796 ms

  Hash Join  [Inner]
    Actual time:   startup=0.012 ms  total=9.563 ms
    Actual rows:   9891.0  |  Hash Cond: (e.dept = t.dept)

    +-  Seq Scan on employee as e  (Outer)
        Actual time:   total=3.595 ms  |  Actual rows: 100000

    +-  Hash  (Inner)  Peak memory: 9 KB
        Actual time:   startup=0.006 ms
        Actual rows:   10.0

        +-  Seq Scan on techdept as t
            Actual rows:   10.0
\end{verbatim}}

\textbf{MariaDB:}

{\small
\parskip0pt\begin{verbatim}
  Optimization time: 0.0633 ms
  Total time:        10.5407 ms

    +-  Table: T  [ALL]
        Key: none  |  Rows: 10  |  pages_accessed=1
        Time: table=0.0067 ms

    +-  Table: E  [ref]  Key: idx_dept_emp
        Ref: mariadb.T.dept
        Loops: planned=10  actual=10
        Rows:  estimated=1  actual=989.1  |  pages_accessed=19829
        Time:  table=9.6074 ms
\end{verbatim}}

Give an interpretation of the execution plan, i.e., describe how the rewritten query is evaluated.

The underlying access pattern is the same as in the original query for both databases: the same operators, indexes, and row counts appear. The differences are in overhead and cache state.

\textbf{PostgreSQL:} The hash join plan is structurally identical. Planning drops to 0.093\,ms (no subquery unnesting needed) and the \texttt{employee} scan drops to 3.595\,ms due to warm buffer cache.

\textbf{MariaDB:} The two-step index-ref access is preserved, but the \textit{Duplicates Removal} node is eliminated because a plain \texttt{JOIN} does not require it. Optimisation time collapses from 10.505\,ms to 0.063\,ms and the \texttt{techdept} scan from 1.658\,ms to 0.007\,ms (warm cache); the \texttt{employee} index lookups still access 19,829 pages (9.607\,ms).

Discuss, how the execution plan changed between the original and the rewritten query. In both the interpretation of the query plans and the discussion focus on the crucial parts, i.e., the parts of the query plans that cause major runtime differences.

\textbf{PostgreSQL:} The physical plan is identical in both cases. The speedup (39.585\,ms $\to$ 9.796\,ms) comes from eliminating subquery-unnesting planning overhead and warm-buffer-cache effects on the \texttt{employee} scan.

\textbf{MariaDB:} The structural change is more significant. The original plan wraps execution in a \textit{Duplicates Removal} node to prevent duplicate employee rows when multiple \texttt{techdept} entries could match. The explicit \texttt{JOIN} allows MariaDB to drop this node entirely. Combined with warm-cache savings on the \texttt{techdept} scan and a 165$\times$ reduction in optimisation time, total execution drops from 66.711\,ms to 10.541\,ms.

\paragraph{Experiment}

Give the runtimes of the original and the rewritten query.

\begin{table}[H]
  \centering
  \begin{tabular}{l|r|r}
    & PostgreSQL [sec] & MariaDB [sec] \tabularnewline
    \hline
    Original query  & 0.040 & 0.067 \tabularnewline
    Rewritten query & 0.010 & 0.011 \tabularnewline
  \end{tabular}
\end{table}

Discuss, why the rewritten query is (or is not) faster than the original query.

The rewritten query is faster in both databases. In \textbf{PostgreSQL} the gain is primarily a planning and warm-cache effect; the physical plan is unchanged. In \textbf{MariaDB} the gain is more structural: dropping the \textit{Duplicates Removal} node and reducing optimisation time by a factor of $\sim$165 are the dominant contributors, with warm cache providing an additional benefit on the \texttt{techdept} scan.

\subsubsection*{Query 2}

\paragraph{Original Query}

Give the second type of query that might be hard for your database to optimize.

\begin{lstlisting}[style=dbtsql]
SELECT E.ssnum 
FROM employee E 
JOIN student S ON E.name = S.name
\end{lstlisting}

\paragraph{Rewritten Query}

Give the rewritten query.

\begin{lstlisting}[style=dbtsql]
SELECT E.ssnum 
FROM employee E 
JOIN student S ON E.ssnum = S.ssnum
\end{lstlisting}

\paragraph{Evaluation of the Execution Plans}

Give the execution plan of the original query.

\textbf{PostgreSQL:}

{\small
\parskip0pt\begin{verbatim}
  Planning time:   3.623 ms
  Execution time:  60.877 ms

  Hash Join  [Inner]
    Actual time:   startup=60.790 ms  total=60.791 ms
    Actual rows:   0.0  |  Hash Cond: (e.name = s.name)

    +-  Seq Scan on employee as e  (Outer)
        Actual time:   total=3.589 ms  |  Actual rows: 100000
        Shared bufs: hit=976  read=0

    +-  Hash  (Inner)  Peak memory: 5408 KB
        Actual time:   startup=37.321 ms  |  Actual rows: 100000

        +-  Seq Scan on student as s
            Actual time:   total=28.407 ms  |  Actual rows: 100000
            Shared bufs: hit=0  read=832
\end{verbatim}}

\textbf{MariaDB:}

{\small
\parskip0pt\begin{verbatim}
  Optimization time: 3.8446 ms
  Total time:        545065.8404 ms

    +-  Table: S  [ALL]
        Key: none  |  Rows: 100000  |  pages_accessed=361
        Time: table=67.0303 ms  other=10.1322 ms

    +-  Block Nested Loop on employee  [BNL]
        access_type: ALL  |  possible_keys: idx_name_emp
        r_loops: 100000  |  Rows: 100000  |  pages_accessed=5954
        Buffer: flat 256 KB
        Time: table=681.7791 ms  unpack=326747.1621 ms  other=217525.0215 ms
        Condition: employee.name = student.name  |  r_filtered: 0%
\end{verbatim}}

Give an interpretation of the execution plan, i.e., describe how the original query is evaluated.

\textbf{PostgreSQL} uses a hash join on the \texttt{TEXT} column \texttt{name}. All 100,000 \texttt{student} rows are read from disk (832 page reads, 28.407\,ms) to build a large hash table (5,408\,KB, 37.321\,ms). All 100,000 \texttt{employee} rows are then scanned (3.589\,ms, cached) and each \texttt{name} probed against the table. Zero rows match (names follow non-overlapping patterns). The wide text key makes the hash table expensive in both memory and build time.

\textbf{MariaDB} uses a Block Nested Loop (BNL) join. First, all 100,000 \texttt{student} rows are scanned (361 pages, 67.030\,ms). Then \texttt{employee} is read in 256\,KB buffer-sized chunks into a flat buffer and joined against the buffered student rows by comparing \texttt{name} strings. Although physically \texttt{employee} is scanned only 13 times (not 100,000), the hash condition \texttt{employee.name = student.name} filters out all rows (r\_filtered = 0\%), so 100,000 unpack iterations yield zero matches. The dominant cost is unpacking and probing the buffer (326,747\,ms), making total execution 545,066\,ms, which is far worse than PostgreSQL, because string comparison inside the BNL loop is extremely expensive at this scale.

Give the execution plan of the rewritten query.

\textbf{PostgreSQL:}

{\small
\parskip0pt\begin{verbatim}
  Planning time:   1.354 ms
  Execution time:  28.853 ms

  Hash Join  [Inner]
    Actual time:   startup=28.775 ms  total=28.776 ms
    Actual rows:   0.0  |  Hash Cond: (e.ssnum = s.ssnum)

    +-  Seq Scan on employee as e  (Outer)
        Actual time:   total=3.572 ms  |  Actual rows: 100000
        Shared bufs: hit=976  read=0

    +-  Hash  (Inner)  Peak memory: 4540 KB
        Actual time:   startup=11.012 ms  |  Actual rows: 100000

        +-  Seq Scan on student as s
            Actual time:   total=4.434 ms  |  Actual rows: 100000
            Shared bufs: hit=832  read=0
\end{verbatim}}

\textbf{MariaDB:}

{\small
\parskip0pt\begin{verbatim}
  Optimization time: 1.6153 ms
  Total time:        51.6597 ms

    +-  Table: E  [index]  Key: idx_ssnum_emp
        Rows: 100000  |  pages_accessed=104
        Time: table=19.1633 ms

    +-  Table: S  [eq_ref]  Key: idx_ssnum_stud
        Ref: mariadb.E.ssnum
        Loops: planned=99854  actual=100000
        Rows:  estimated=1  actual=0  |  pages_accessed=200000
        Time:  table=26.6500 ms
\end{verbatim}}

Give an interpretation of the execution plan, i.e., describe how the rewritten query is evaluated.

\textbf{PostgreSQL} switches the join key to the 4-byte integer \texttt{ssnum}. The \texttt{student} hash table shrinks to 4,540\,KB and is built in 11.012\,ms; its pages are now cached from the previous run (0 disk reads, 4.434\,ms scan). The \texttt{employee} scan is unchanged. Zero rows match because \texttt{ssnum} ranges do not overlap.

\textbf{MariaDB} chooses a fundamentally different plan. It scans \texttt{employee} via \texttt{idx\_ssnum\_emp} (104 pages, 19.163\,ms), then for each of the 100,000 rows performs an \texttt{eq\_ref} lookup into \texttt{student} using \texttt{idx\_ssnum\_stud} (200,000 index page accesses, 26.650\,ms). This index nested-loop join on integers is vastly cheaper than the previous BNL plan, since integer point lookups via an index are far less expensive than buffer-based string matching across 100,000 iterations.

Discuss, how the execution plan changed between the original and the rewritten query. In both the interpretation of the query plans and the discussion focus on the crucial parts, i.e., the parts of the query plans that cause major runtime differences.

\textbf{PostgreSQL:} The operator type (hash join over full table scans) is unchanged. Switching from \texttt{TEXT} to \texttt{INT} key reduces the hash table by 868\,KB and cuts build time from 37.321\,ms to 11.012\,ms. The remaining speedup comes from the warm cache eliminating 832 disk reads for the \texttt{student} scan (28.407\,ms $\to$ 4.434\,ms).

\textbf{MariaDB:} The plan changes fundamentally. The original BNL join cannot exploit any index on the \texttt{name} join condition and spends 326,747\,ms unpacking and comparing string buffers, yielding a total of 545,066\,ms. Joining on \texttt{ssnum} allows MariaDB to replace the BNL entirely with an index nested-loop join using \texttt{idx\_ssnum\_emp} and \texttt{idx\_ssnum\_stud}, reducing total time to 51.660\,ms, an improvement of roughly $10{,}500\times$. The key insight is that BNL with a string predicate is catastrophically expensive, while an \texttt{eq\_ref} index lookup on an integer key is cheap even for 100,000 iterations.

\paragraph{Experiment}

Give the runtimes of the original and the rewritten query.

\begin{table}[H]
  \centering
  \begin{tabular}{l|r|r}
    & PostgreSQL [sec] & MariaDB [sec] \tabularnewline
    \hline
    Original query  & 0.061 & 404.151 \tabularnewline
    Rewritten query & 0.029 & 0.052 \tabularnewline
  \end{tabular}
\end{table}

Discuss, why the rewritten query is (or is not) faster than the original query.

In \textbf{PostgreSQL} the improvement ($\sim$$2\times$) comes from a smaller, faster hash table for integer keys and warm-cache effects on the \texttt{student} scan. The plan structure is unchanged.

In \textbf{MariaDB} the improvement is dramatic (∼\sim
∼$10{,}500\times$). Joining on \texttt{name} forces a Block Nested Loop join with no usable index: the 256\,KB buffer must be unpacked and compared against 100,000 student rows per iteration, with string matching consuming 326,747\,ms in unpack time alone. Joining on \texttt{ssnum} replaces the BNL entirely with an index nested-loop join via \texttt{idx\_ssnum\_emp} and \texttt{idx\_ssnum\_stud}. This shows that the choice of join attribute matters far more for MariaDB than for PostgreSQL in this scenario: PostgreSQL's hash join handles text keys reasonably well, while MariaDB falls back to an extremely expensive BNL when no index-friendly join key is available.

\subsection*{Time Spent on this Assignment}

Time in hours per person:
\begin{itemize}
    \item Kitajima Mitsuhiro: \textbf{5,0}
    \item Sadykova Nargiz: \textbf{2,5}
    \item El Lbidi Omar: \textbf{1,5}
\end{itemize}

\subsection*{References}

\begin{table}[H]
  \centering
  \begin{tabular}{c}
    \hline
    \textbf{Important:} Reference your information sources! \tabularnewline
    Remove this section if you use footnotes to reference your information sources. \tabularnewline
    \hline
  \end{tabular}
\end{table}

\end{document}